\chapter{Цепная степень как родитель модульных боровских частот}
\label{ch:v8-modular-bohr-parent-origin-gate}

\section{Поиск среди уже выведенных операторов}

Предыдущая глава показала, что нетривиальное KMS-состояние требует
гамильтониана, относительно которого межугловые стрелки имеют определённые
боровские частоты. Проверим четыре внутренних кандидата:
\begin{enumerate}
 \item полярную пару Грамма $A_0^*A_0\oplus A_0A_0^*$;
 \item ограничение на концевые секторы Ходж-метрики;
 \item семейные Гиббс-состояния Тома IV;
 \item оператор степени связывающей цепи Тома VII.
\end{enumerate}

Первые три кандидата не закрывают задачу. Пара Грамма сплетает фоновую
инцидентность с нулевой частотой, но не даёт общей частоты межсекторных стрелок
семейств: минимальные относительные остатки равны $1/\sqrt2$ для $QLYR$ и
$1$ для $XLdR$. Ходж-кандидат имеет разрыв $\log\eta$ с невыведенным
$\eta$, а семейные состояния живут в $M_3$, не имея канонического подъёма
в $M_{11}\oplus M_{10}$.

\section{Канонический целочисленный кандидат}

Связывающий комплекс уже несёт оператор степени
\begin{equation}
 N=\operatorname{diag}(0I_{11},I_{21},2I_{10}).
 \label{eq:v8-chain-number-full}
\end{equation}
На концевой алгебре остаётся
\begin{equation}
 N_{\partial}=0I_{11}\oplus2I_{10}.
 \label{eq:v8-chain-number-endpoint}
\end{equation}
Для любой из тринадцати связывающих и межсекторных матриц стрелок
$V:\mathbb C^{11}\to\mathbb C^{10}$ выполняется точное тождество
\begin{equation}
 [N_{\partial},V]=2V,
 \qquad
 [N_{\partial},V^*]=-2V^*.
 \label{eq:v8-chain-number-bohr-pair}
\end{equation}
Таким образом, один ранее выведенный оператор одновременно типизирует все
межугловые стрелки и не содержит непрерывного параметра.

При каноническом безразмерном Гиббс-представителе
\begin{equation}
 \rho_N=\frac{e^{-N_{\partial}}}{\Tr e^{-N_{\partial}}}
 =aI_{11}\oplus bI_{10},
 \qquad
 \frac ba=e^{-2},
 \label{eq:v8-chain-number-gibbs-state}
\end{equation}
направленные скорости удовлетворяют
\begin{equation}
 \frac{\gamma_\uparrow}{\gamma_\downarrow}=e^{-2}.
 \label{eq:v8-chain-number-rate-ratio}
\end{equation}

\section{Полный направленный процесс}

Все самосопряжённые скачки переноса заменены парами $V,V^*$, а калибровочные
слагаемые оставлены без изменения. Полученный генератор является
унитальным в картине Гейзенберга, вполне положительным и KMS-симметричным
относительно $\rho_N$. Машинный аудит даёт
\begin{align}
 \dim\ker\mathcal L_N&=1,
 &\lambda_{\rm gap}&=0.0212599674\ldots,\nonumber\\
 \Pr(\mathbb C^{11})&=0.8904465168\ldots,
 &\Pr(\mathbb C^{10})&=0.1095534832\ldots.
 \label{eq:v8-chain-number-forward-process}
\end{align}
Остатки KMS-симметрии и стационарности меньше $10^{-16}$. При нулевом
разрыве конструкция точно возвращает прежний симметричный генератор.

Это первый положительный внутренний источник формы направленных скоростей.
Однако он пока задаёт безразмерную модульную динамику, а не физическую
температуру или секунду.

\section{Оставшаяся дискретная развилка}

В Томе VII оператор степени был определён с точностью до обращения цепи
\begin{equation}
 N\longmapsto2I-N.
 \label{eq:v8-chain-number-reversal}
\end{equation}
Для нормы относительной кривизны это обращение меняло только знак коммутатора и
потому было невидимо. Для KMS-состояния оно существенно:
\begin{equation}
 \frac{b_{\rm rev}}{a_{\rm rev}}=e^2.
\end{equation}
Обращённый процесс также вполне положителен, KMS-симметричен и примитивен;
его щель равна $0.1765195742\ldots$. Поэтому текущий родитель выделяет не
один процесс, а дискретную пару. Направление коизометрии
$\mathbb C^{11}\to\mathbb C^{10}$ и индексный дефект ранга один дают
естественную зацепку, но прежнее действие не доказало, что именно они
выбирают знак $N$.

\begin{theorem}[Условный цепно-модульный родитель]
Оператор степени связывающего комплекса является корректно типизированным,
калибровочно-инвариантным и не содержащим непрерывного параметра общим боровским
гамильтонианом всех стрелок переноса Тома VIII. После направленного
разбиения он порождает примитивную KMS-полугруппу с отношением скоростей
$e^{-2}$. Результат условен относительно выбора ориентации цепи и
безразмерной Гиббс-нормировки.
\end{theorem}

\begin{nogocriterion}[Запрет преждевременного чтения $e^{-2}$ как числа природы]
Нельзя объявлять $e^{-2}$ физическим отношением скоростей, пока действие
родителя не выбрало $N$ против $2I-N$ и не связало целочисленную степень с
физической энергией и временем. Обращённая ветвь $e^2$ проходит те же
операторные тесты.
\end{nogocriterion}

\begin{protocol}
\begin{enumerate}
 \item корректно типизированный концевой гамильтониан: \textbf{получен};
 \item источник: \textbf{оператор степени цепи Тома VII};
 \item общий разрыв всех стрелок переноса: \textbf{$2$};
 \item непрерывные параметры в разрыве: \textbf{нет};
 \item направленный Линдблад-процесс: \textbf{построен};
 \item KMS-симметрия: \textbf{пройдена};
 \item примитивность: \textbf{сохранена};
 \item прямое отношение скоростей: \textbf{$e^{-2}$};
 \item обращённое отношение скоростей: \textbf{$e^2$};
 \item ориентация выбрана прежним действием: \textbf{нет};
 \item физическая единица энергии и времени: \textbf{не выведена};
 \item следующий тест: \textbf{индексно-дефектный селектор ориентации цепи}.
\end{enumerate}
\end{protocol}

Машинный сертификат сохранён в
\path{s2t_v8_modular_bohr_parent_origin_gate_results.json}.

\begin{thebibliography}{9}
\bibitem{v8-bohr-ding}
Z.~Ding, B.~Li, L.~Lin,
\emph{Efficient quantum Gibbs samplers with Kubo--Martin--Schwinger
detailed balance condition}, arXiv:2404.05998.
\bibitem{v8-bohr-fagnola}
F.~Fagnola, V.~Umanit\`a,
\emph{Generators of detailed balance quantum Markov semigroups},
arXiv:0707.2147.
\bibitem{v8-bohr-project}
Том VII, гейт цепной степени и относительное факторпространство конуса отображения.
\end{thebibliography}